\documentclass[11pt]{article}
\usepackage[a4paper,margin=1in]{geometry}
\usepackage[T1]{fontenc}
\usepackage{lmodern,microtype}
\usepackage{amsmath,amssymb,amsthm,mathtools}
\usepackage{booktabs}
\usepackage{xcolor}
\usepackage[numbers,sort&compress]{natbib}
\usepackage[colorlinks=true,linkcolor=blue!55!black,citecolor=blue!55!black]{hyperref}

\newtheorem{theorem}{Theorem}[section]
\newtheorem{proposition}[theorem]{Proposition}
\newtheorem{corollary}[theorem]{Corollary}
\theoremstyle{remark}
\newtheorem{remark}[theorem]{Remark}

\title{A Certified Deterministic-Numerator Boundary Budget}
\author{Bin Wang}
\date{July 2026}

\begin{document}
\maketitle

\begin{abstract}
We turn the all-order Cauchy interface of RH-252 into a usable numerical
certificate.  The exact RH-15 factorization and the RH-13 Arb bounds imply,
without angular sampling, that the normalized Hardy-scaled deterministic
numerator satisfies $M_{7/5}<108$.  The order-28 anchor of RH-253 then gives
an order-29 unit-disk logarithmic tail below $0.021866475$ and relative
multiplicative error below $0.022107298$.  This certifies only the target
boundary constant.  It supplies neither a legal cloud head, a cloud-to-target
coefficient bridge, nor a uniform quotient tail.
\end{abstract}

\section{Factorization and nuclear input}

Let $G$ be the one-step deterministic numerator and let
$G_H(z)=G(z/r_H)$ with $r_H=17/20$.  RH-15 proves the exact factorization
\citep{WangRH15}
\begin{equation}
 G(u)=e^{P_1u}\widetilde D_{1,+}(u^2)A_*(u^2)B(u^2)C(u),
 \qquad \widetilde D_{1,+}(w)=\det(I-wT).
 \label{eq:factor}
\end{equation}
The logarithms below are the branches normalized to vanish at the origin.
RH-13 realizes $T$ on a Wiener algebra of radius $R=0.7$ and certifies its
finite section and tail with Arb arithmetic \citep{WangRH13}.

For $k\ge1$, the nonzero reduced columns have the form
\[
 C_{2k}=2a(x)\left(t(x)^k/R^{2k}-m_{2k}\right),
 \qquad
 m_{2k}=\binom{2k}{k}\left(\frac{r}{2R}\right)^{2k}.
\]
If $\tau=\|t\|_R/R^2<1$, summing the column norms gives the nuclear bound
\begin{equation}
 \nu_1:=\nu(T)\le2\|a\|_R\left\{
 \frac{\tau}{1-\tau}
 +\frac1{\sqrt{1-(r/R)^2}}-1\right\}.
 \label{eq:nuclear}
\end{equation}
Indeed, the first term sums $\tau^k$, while the second follows from
$\sum_{k\ge0}\binom{2k}{k}x^k=(1-4x)^{-1/2}$.

Let $M$ be the RH-13 finite section and let $\varepsilon$ bound the operator
remainder.  Then
\begin{equation}
 p_1=\|M\|+\varepsilon\ge\|T\|,
 \qquad
 p_2=\|M^2\|+2\|M\|\varepsilon+\varepsilon^2\ge\|T^2\|.
 \label{eq:p12}
\end{equation}
Thus $\nu(T^2)\le p_1\nu_1$ and $\nu(T^3)\le p_2\nu_1$ by the nuclear-ideal
inequality.

\section{A uniform boundary theorem}

Put $S=7/5$, $U=S/r_H=28/17$, $W=U^2$, and $t=U/\lambda$.
The choice is strict: $U<\lambda$.  Write
$q\ge\|T^3\|$ for the RH-13 certified cube bound.

\begin{theorem}[Modulo-three Fredholm majorant]\label{thm:fredholm}
If $qW^3<1$, then throughout $|w|=W$,
\begin{equation}
 |\log\det(I-wT)|
 \le
 \frac{\nu_1W+(p_1\nu_1)W^2/2+(p_2\nu_1)W^3/3}
 {1-qW^3}.
 \label{eq:fredholm}
\end{equation}
\end{theorem}

\begin{proof}
For $j=1,2,3$ and $k\ge0$, the nuclear-ideal inequality gives
$\nu(T^{3k+j})\le q^k\nu(T^j)$.  Combine this with
$|\operatorname{tr}T^n|\le\nu(T^n)$ in the Fredholm logarithm, group the
series modulo three, and replace $3k+j$ in the denominator by $j$.
The remaining series is geometric with ratio $qW^3$.
\end{proof}

The endpoint factors also admit exact scalar bounds.  With
$b_n=\lambda^{-n}/(1+\lambda^{-n})$ and
$d_n=\lambda^{-2n}/(1-\lambda^{-2n})$, coefficientwise absolute majorants
give
\begin{align}
 |\log A_*(u^2)|&\le-\log(1-t^2),\label{eq:A}\\
 |\log B(u^2)|&\le
 \frac{-\log(1-t^2)}{2(1-\lambda^{-2})}.\label{eq:B}
\end{align}
The identity $P_1=b_1$ cancels the linear term of $\log C$, so
\begin{equation}
 |P_1u+\log C(u)|\le\operatorname{artanh}(t)-t.
 \label{eq:C}
\end{equation}

\begin{theorem}[Certified boundary budget]\label{thm:budget}
For the normalized logarithm of $G_H$,
\[
 \boxed{M_{7/5}:=\sup_{|z|=7/5}|\log G_H(z)|<107.906078<108.}
\]
\end{theorem}

\begin{proof}
The 200-decimal-place Arb replay of the RH-13 certificate gives
\begin{center}
\small
\begin{tabular}{lr@{\qquad}lr}
\toprule
$\nu_1$ & $<4.623248864$ & $p_1$ & $<0.633964866$\\
$p_2$ & $<0.174350001$ & $qW^3$ & $<0.715126024$\\
Fredholm term & $<100.715071$ & $A_*$ term & $<3.291531$\\
$B$ term & $<2.551222$ & linear+$C$ term & $<1.348256$\\
\bottomrule
\end{tabular}
\end{center}
All entries are outward endpoints, and every displayed strict comparison is
checked again at 100 and 150 decimal places.  Summing
\eqref{eq:fredholm}--\eqref{eq:C} proves the claim.  This is a uniform
factor majorant on the entire circle, not a sampled maximum.
\end{proof}

\section{Order-29 consequence and boundary}

RH-252 proves that Cauchy's estimate converts $M_S$ into an all-order target
tail \citep{WangRH252}; RH-253 supplies deterministic coefficients through
order 28 \citep{WangRH253}.  Hence the first omitted order is $N=29$.

\begin{corollary}[Certified unit-disk target tail]
On $|z|\le1$,
\[
 \sum_{n\ge29}\frac{|a_n|}{n}
 \le M_{7/5}\frac{(5/7)^{29}}{1-5/7}
 <0.021866475.
\]
The corresponding relative multiplicative error is below
$e^{0.021866475}-1<0.022107298$.
\end{corollary}

The protocol reruns the RH-13 Arb certificate at 100, 150, and 200 decimal
places, evaluates only interval expressions, checks twelve outward decimal
claims, and archives source hashes.  It deliberately does not use the
mislabelled $\lambda^2$ radius field in an older RH-46 JSON: the one-step
zero-free radius used here is $\lambda$.

In the five-component RH-260 ledger \citep{WangRH260}, the status changes
from $(0,0,0,1,0)$ to $(0,0,0,1,1)$.  The legal anchored head, current-cloud
coefficient bridge, and uniform quotient tail remain open, so the complete
certificate count is zero.  A finite order-28 anchor is not promoted to an
all-order identification.  Gates A--E remain false/open.  No Hilbert--Polya
operator, Riemann-zero identification, completed-zeta divisor equality, or
implication of RH is asserted.

\bibliographystyle{plainnat}
\bibliography{references}
\end{document}
